\section{Evaluation Method}\label{sec:experimental}
We ask whether temporal history changes migration decisions, what each
admission component contributes, how destinations are distributed, and how
much computation the decision procedure requires. We compare the complete
procedure with component ablations under the same trace-replay protocol.

\subsection{Trace and event workload}
The study uses Alibaba Microservices Trace 2021
\citep{luo2021alibaba,alibaba2021trace}, combining \texttt{Node\_0} with
\texttt{MSResource\_0}. The full node archive contains 18,868,400 records,
13,118 nodes, and 1,440 timestamps at 30-second intervals. Removing
incomplete full trajectories gives 13,058 nodes. Local alignment with the
clean microservice shard yields 12,727 nodes over 120 timestamps and
1,527,240 node observations. The microservice shard contains 11,561,072
clean records, 1,303 services, and 96,346 service--instance pairs.

Instance CPU signals are used to trigger decisions, not to infer node
capacity. After a 300-second warm-up, an event is emitted when the instance
CPU signal increases by more than 0.05 over one interval. The rule produces
20,419 events; excluding the last timestamp, whose outcome is unavailable,
leaves 20,324 events over 108 timestamps. This is a constructed workload
of placement reconsiderations, not a trace of new service arrivals.

The natural workload preserves the observed service-frequency imbalance:
its ten most frequent services account for approximately 46\% of events.
A second workload retains at most two events per service per timestamp,
using deterministic selection. It contains 6,875 events. The balanced test
subset has a top-ten service share of 22.55\%. Both test workloads use the
same 59 timestamps, so they are sensitivity views of the same trace rather
than independent replications.

\subsection{Information boundary and model selection}
An event at $t$ uses the instance change ending at $t$. Its decision uses
node observations at $t,t-\Delta,t-2\Delta$; evaluation reads the selected
node's recorded state at $t+\Delta$. The replay leaves that future state
unchanged. A later event still refers to the source recorded in the trace;
the evaluation is not a closed-loop simulation of the resulting placements.

The chronological split is at 1,800,000\,ms. The reported partitions are
shown in Table~\ref{tab:partitions}. Model fitting uses development data;
estimator selection and operating-point selection use the pre-test region.
The final test region is reserved for evaluation. Exact split inequalities,
label-boundary handling, and the provenance of projection features used
during fitting must be recovered from the experiment code to audit this
information boundary fully.

\begin{table}[t]
\centering
\caption{Reported chronological partitions.}\label{tab:partitions}
\begin{tabular}{lrr}
\toprule
Partition & Transitions & Natural events\\
\midrule
Development &432,718&--\\
Validation &190,905&--\\
Pre-test total &623,623&8,695\\
Test &750,893&11,629\\
\bottomrule
\end{tabular}
\end{table}

The balanced pre-test and test subsets contain 3,088 and 3,787 events.
The principal model is scikit-learn histogram gradient boosting
\citep{pedregosa2011sklearn}. Its direct predictive control uses current
telemetry without temporal lags or changes. Both use the same chronological
partitions. Reported model metrics include ROC-AUC, PR-AUC, Brier score,
log loss, and expected calibration error. A separate probability-calibration
transform is not specified; a low calibration error alone does not identify
such a transform.

The temporal threshold is $0.011362695354780319$; the current-state
threshold is $0.004292615073088162$. These were selected at matched
pre-test operating points. A stricter temporal threshold,
$0.0020594307448572016$, is used for sensitivity analysis. The empirical
headroom margins are
\begin{align}
q_c&=0.025929729975254,\\
q_m&=0.00021962105567108203.
\end{align}
The operating boundary is 0.8 for both resources. Deterministic headroom
cutoffs of 1, 0.99, and 0.975 provide comparison points. The default
shortlist size is $k=256$, compared with 128, 512, 1024, 2048, and full
ranking. Reservation capacity is $C=1$, with $C=2$ as a sensitivity check.
Settings were selected before final testing; the selection logs are not
included with this draft.

\subsection{Synthetic policy attributes}\label{sec:policyoverlay}
The trace has no jurisdiction, trust, encryption, or isolation labels.
The evaluation assigns them using a stable hash of identifiers and a policy
seed, independently of resource telemetry. Nodes receive one of four
jurisdictions, one of three trust ranks, and Boolean encryption and
trusted-execution attributes. Services receive one of six policy classes.
Under the reference seed, legal-node counts are 12,727 (global public),
10,320 (global protected), 6,206 (EU restricted), 2,613 (US restricted),
1,501 (APAC restricted), and 866 (EU confidential).

Thirty deterministic overlays vary node capabilities and service
requirements while retaining the trace and model. This tests behavior
under the chosen policy generator, not compliance with actual Alibaba
policies. In particular, the native source baseline can violate an overlay
because that overlay was assigned after collection. These violations do
not describe a production compliance failure.

\subsection{Controls and ablations}
All methods share events, event order, telemetry, and fixed tie breaking.
Applicable methods also share the policy overlay, active-source exclusion,
and reservation rules. Each method retains its source if that source passes
its own admission rule. Consequently, comparisons measure both retention
and destination choice.

\begin{itemize}[leftmargin=*]
\item \emph{Observed source} retains the trace source without admission.
\item \emph{Random-policy-Q95} and \emph{Least-loaded-policy-Q95} enforce
policy and headroom, then choose a random or minimally utilized node.
Random selection is repeated with 20 deterministic seeds.
\item \emph{Policy-only-fair} uses policy and selection history.
\emph{Policy-Q95-fair} adds headroom; \emph{Policy-boundary-0.99} uses a
stricter headroom cutoff.
\item \emph{Policy-current-HGB-Q95} replaces temporal features with
current-state learned risk.
\item \emph{Temporal-Q95-no-policy} removes the policy gate;
\emph{Policy-temporal-no-projection} removes the headroom check.
\end{itemize}
Together with \cage, these are ten methods, including the observed-source
control. They are internal controls, not implementations of published
orchestrators. No direct performance comparison with Kubernetes, NetMARKS,
or another deployed scheduler is claimed.

\subsection{Metrics and uncertainty}
Migration rate counts decisions whose selected node differs from the
recorded source. Policy, temporal, and Q95 violation rates count failures
of $L(s,n)=1$, $p_{n,t}\leq\theta$, and $g_{n,t}\leq1$, respectively.
The temporal metric applies \cage's temporal screen to all methods,
including those optimized for another admission rule. It measures agreement
with that screen and is not an independent measure of predictive quality.

The native next-state unsafe rate averages $y_{n,t+\Delta}$ for selected
nodes. It describes their unmodified recorded trajectories. It does not
measure overload caused or prevented by migration. Search expansion and
infeasible-event counts are reported separately; no infeasible decisions
are reported in the evaluated configurations.

For migrated-destination counts $z_n$, Jain's index is
\begin{equation}
J=\frac{(\sum_n z_n)^2}{M\sum_n z_n^2}.
\end{equation}
It includes nodes with zero assignments and measures destination spread,
not fairness among tenants or services. For the no-migration control it is
undefined; the tables mark it accordingly. Top-ten share is the fraction
of migrations assigned to the ten most-used nodes.

The paired timestamp-cluster bootstrap samples the 59 test timestamps
with replacement and includes every event in each sampled timestamp.
The same samples are applied to both methods. Random-baseline replicate
identity is also sampled where applicable. Reported intervals are
percentile 95\% intervals from 20,000 resamples of the difference
\cage minus baseline. This preserves within-timestamp dependence but
does not preserve correlation between adjacent timestamps or propagate
uncertainty from fitting a different risk model.

\subsection{Timing and artifact status}
One latency sample is recorded per test timestamp. The in-memory path
includes feature construction, all required model calls, filtering,
shortlisting, ranking, reservations, and final checks. Telemetry collection,
network and file I/O, serialization, binding, and migration execution are
excluded. The reported platform is a Dell Precision 3561 running Ubuntu
under WSL2 with Python 3.12.3. The exact CPU, memory allocation, thread
settings, warm-up protocol, and repeat count are not specified, limiting
reproducibility of the absolute timing results.

The available package contains manuscript sources and rendered figures.
It does not contain the preprocessing code, model, policy generator,
replay implementation, environment lockfile, or raw measurements needed
to reproduce these results. Artifact availability should therefore not be
inferred from the public availability of the source trace.
